\chapter{Cubic curves}\label{ch:cubics}

\begin{center}\itshape
Every curve of the third order is the shadow of one of five parabolas.
\upshape(paraphrase of Newton, \emph{Enumeratio linearum tertii ordinis}, 1704)
\end{center}

Around 1621 Claude-Gaspard Bachet de M\'eziriac, whose Latin edition of Diophantus' \emph{Arithmetica} came out that year (it is in a copy of this
book that Fermat later wrote his famous marginal note), looked at the equation $y^2=x^3-2$. It has the obvious solution $x=3$, $y=5$. Is there another
one in rational numbers? Bachet found one, $x=\frac{129}{100}$, $y=\frac{383}{1000}$, and a recipe to produce more: draw the tangent to the curve at
$(3,5)$ and see where it cuts the curve again (Figure~\ref{fig:bachet}). Today we see in his trick the first step of a group law: on a cubic curve, two
points can be ``added'' to give a third, and the rational points form a group.

Newton, who classified the real cubic curves into $72$ species in an appendix to his \emph{Opticks} (1704), noticed that all of them are projections of
the five ``divergent parabolas'' $y^2=ax^3+bx^2+cx+d$. In this chapter we prove a sharp form of his remark, the Weierstrass normal form, and then follow
the chord and the tangent: the group law, its associativity from the Cayley--Bacharach theorem of Chapter~7, the nine inflection points and their twelve
lines, which form the finite plane of order $3$ of Chapter~4, and the curves over finite fields, which a cryptographer would recognise.

\begin{figure}[ht]
\centering
\fig{ch08-bachet}
\caption{Bachet's tangent on $y^2=x^3-2$. The tangent at $P=(3,5)$, of slope $\frac{27}{10}$, meets the curve again at
$P*P=\bigl(\frac{129}{100},\frac{383}{1000}\bigr)$; reflecting in the $x$-axis gives the point $2P$ of Section~\ref{sec:weierstrass}.}
\label{fig:bachet}
\end{figure}

\begin{think}
A line meets a cubic in three points. If two of them are rational, why must the third be rational too? What goes wrong with conics, where a line meets
the curve in two points?
\end{think}

\section{Chords and tangents}

Throughout, $K$ is a field and $E$ is a \emph{smooth} cubic curve over $K$, given by a homogeneous cubic $F$ with no singular point over an algebraic
closure $\overline K$. A smooth cubic is irreducible (the solution of Exercise~\ref{ex:polarcurve}), so no line is contained in $E$. We write $E(K)$ for the
set of points of $E$ with coordinates in $K$.

Let $L$ be a line through two points $[p]$ and $[v]$. Its points are $[sp+tv]$, and $F(sp+tv)$ is a binary form of degree $3$ in $(s,t)$, not zero since $L$ is
not a component of $E$. Its roots are the points of $L\cap E$; their \emph{multiplicities} add up to $3$ over $\overline K$.

\begin{lemma}\label{lem:tangentdouble}
Let $P=[p]$ be a point of $E$ and $L$ a line through $P$. The root $P$ of $F$ restricted to $L$ is at least double if and only if $L$ is the tangent line
at $P$.
\end{lemma}

\begin{proof}
Write $L$ as the points $[p+tv]$ together with $[v]$, where $[v]\ne P$. By Taylor's formula $F(p+tv)=F(p)+t\,\nabla F(p)\cdot v+t^2(\dots)$, and $F(p)=0$.
So $t=0$ is at least a double root if and only if $\nabla F(p)\cdot v=0$, that is, if and only if $[v]$ lies on the tangent line $\nabla F(p)\cdot x=0$.
That line also contains $P$, since $\nabla F(p)\cdot p=3F(p)=0$ by Euler's formula, so this says exactly that $L$ is the tangent.
\end{proof}

For $P,Q\in E(K)$, let $L$ be the line $PQ$, or the tangent at $P$ if $P=Q$. By the lemma, $P$ and $Q$ account for two of the three roots of the binary cubic.
A cubic with coefficients in $K$ and two roots in $K$ has its third root in $K$ too: divide by the two known linear factors. Call the third point $P*Q$; it
lies in $E(K)$. The operation $*$ is commutative, and $P*Q=R$ means that $P,Q,R$ are the three points of a line with $E$, counted with multiplicity; so
$P*(P*Q)=Q$.

The operation $*$ is not a group law: it has no neutral element (Exercise~\ref{ex:noneutral}). The fix is to choose a point.

\begin{definition}
Fix $O\in E(K)$. For $P,Q\in E(K)$, set $P+Q=O*(P*Q)$.
\end{definition}

\begin{proposition}\label{prop:groupbasic}
The operation $+$ is commutative, $O$ is a neutral element, and every $P$ has an inverse $-P=P*(O*O)$.
\end{proposition}

\begin{proof}
Commutativity comes from that of $*$. Next, $P+O=O*(P*O)$, and $O$, $P$, $P*O$ are the three points of one line, so $O*(P*O)=P$. Finally, let $O'=O*O$, the
third point of the tangent at $O$, and $P'=P*O'$. Then $P$, $O'$, $P'$ lie on a line, so $P*P'=O'$, and $P+P'=O*O'=O$, because $O,O,O'$ are the three points
of the tangent at $O$.
\end{proof}

The hard part is associativity. Here is where Chapter~7 pays off.

\begin{theorem}[associativity]\label{thm:assoc}
Let $P,Q,R\in E(K)$ and set
\[
S'=P*Q,\quad S=P+Q,\quad T'=S*R,\qquad U'=Q*R,\quad U=Q+R,\quad T''=P*U.
\]
If the nine points $P,Q,R,O,S,S',U,U',T'$ are distinct, then $(P+Q)+R=P+(Q+R)$.
\end{theorem}

\begin{proof}
Since $(P+Q)+R=O*T'$ and $P+(Q+R)=O*T''$, it suffices to show $T'=T''$. Consider the six lines
\[
\ell_1=PQS',\quad \ell_2=SRT',\quad \ell_3=U'OU,\qquad m_1=S'OS,\quad m_2=QRU',\quad m_3=PUT'',
\]
each named by its three points on $E$ (for instance $U=O*U'$ means that $U'$, $O$, $U$ are collinear). The cubic $C_1=\ell_1\ell_2\ell_3$ meets $E$ in the nine
distinct points $P,Q,S',S,R,T',U',O,U$, and $E$, being irreducible, has no common component with it. The cubic $C_2=m_1m_2m_3$ passes through eight of these
points, $S',O,S,Q,R,U',P,U$. By the Cayley--Bacharach theorem (Theorem~\ref{thm:cayleybacharach}), it also passes through $T'$. So $T'$ lies on $m_1$, $m_2$
or $m_3$, and also on $E$. The points of $E$ on $m_1$ are $S',O,S$, and on $m_2$ they are $Q,R,U'$, all different from $T'$. So $T'$ is on $m_3$, whose points on
$E$ are $P,U,T''$; since $T'\ne P,U$, we get $T'=T''$.
\end{proof}

\begin{figure}[ht]
\centering
\fig{ch08-real}
\caption{Two real cubics: $y^2=x^3-x$ has two components, $y^2=x^3-x+1$ has one. The red points are the real inflection points in the affine plane; with
the point at infinity $O$ they are the three points of the vertical dashed line (Theorem~\ref{thm:realflexes}).}
\label{fig:realcubics}
\end{figure}

\begin{remark}
The theorem covers ``most'' triples. The remaining cases, when some of the nine points coincide, are handled in Tome~12 (Chapter~12, the theorem on the
group law), where Cayley--Bacharach is proved for divisors, that is, for points counted with multiplicity. Over $\mathbb R$ or $\mathbb C$ one can also
argue by continuity: both sides of the identity are continuous functions of $(P,Q,R)$, and triples in general position are dense. The companion notebook
checks associativity for every triple of points on several curves over finite fields, degenerate or not. From now on we use that $(E(K),+)$ is an abelian
group.
\end{remark}

\section{The Weierstrass form}\label{sec:weierstrass}

A point $O$ of $E$ is an \emph{inflection point}, or \emph{flex}, if its tangent line meets $E$ only at $O$, with multiplicity $3$; equivalently, $O*O=O$.
Choosing a flex as the origin makes everything simpler: then $-P=P*O$, and three points add up to $O$ exactly when they are collinear (Exercise~\ref{ex:collinear}).

\begin{proposition}[Weierstrass form]\label{prop:weierstrass}
Let $K$ have characteristic different from $2$ and $3$, and let $E$ be a smooth cubic over $K$ with a flex $O\in E(K)$. In suitable coordinates over
$K$, $O=[0:1:0]$ and
\[
E:\ y^2z=x^3+axz^2+bz^3,\qquad\text{that is}\qquad y^2=x^3+ax+b\ \text{in the chart }z=1,
\]
with $4a^3+27b^2\ne0$. Conversely, every such equation defines a smooth cubic with a flex at $[0:1:0]$.
\end{proposition}

\begin{proof}
By Proposition~\ref{prop:frames} we may send $O$ and a second point of its tangent to $[0:1:0]$ and $[1:0:0]$; then $O=[0:1:0]$ and its tangent is
the line $z=0$. The restriction $F(x,y,0)$ is a binary cubic with a triple root at $[0:1]$, so $F(x,y,0)=cx^3$, and $c\ne0$ since the line $z=0$
is not a component. Hence $F=cx^3+zQ(x,y,z)$ with $Q$ a quadratic form. Let $e$ be the coefficient of $y^2$ in $Q$. If $e=0$, then $F$ has no monomial
$y^3$, $xy^2$ or $y^2z$, so all three partial derivatives vanish at $[0:1:0]$, and $O$ would be singular. So $e\ne0$, and $Q=ey^2+y\,(\alpha x+\beta z)+q(x,z)$.
Replacing $y$ by $y-\frac{\alpha x+\beta z}{2e}$ removes the middle term, and the equation takes the form
\[
ey^2z=c'x^3+c_2x^2z+c_1xz^2+c_0z^3,\qquad c'=-c\ne0.
\]
Replacing $z$ by $\frac{c'}e\,z$ and dividing by $c'$ makes both leading coefficients $1$, and then a substitution $x\mapsto x-\frac s3z$, where $s$ is the
new coefficient of $x^2z$, removes the $x^2z$ term. None of these
substitutions moves $[0:1:0]$, and the equation is now $y^2z=x^3+axz^2+bz^3$.

For the smoothness condition, let $f(x)=x^3+ax+b$. In the chart $z=1$ the equation is $y^2-f(x)=0$, and a singular point has $2y=0$ and $f'(x)=0$, with
$f(x)=y^2=0$: a repeated root of $f$. At $[0:1:0]$ the derivative with respect to $z$ is $y^2-\dots=1\ne0$. A common root of $f$ and $f'=3x^2+a$ satisfies
$x^2=-\frac a3$ and $0=f(x)=x(x^2+a)+b=\frac{2a}3x+b$. If $a=0$ this forces $b=0$. If $a\ne0$ it forces $x=-\frac{3b}{2a}$, and then $x^2=-\frac a3$ reads
$27b^2=-4a^3$; conversely, if $4a^3+27b^2=0$ and $a\ne0$, then $x=-\frac{3b}{2a}$ is a common root. So $E$ is smooth exactly when $4a^3+27b^2\ne0$. Finally
$F(x,y,0)=-x^3$, so the line $z=0$ meets $E$ only at $[0:1:0]$, three times: it is a flex.
\end{proof}

On a Weierstrass curve, the line through $O=[0:1:0]$ and an affine point $(x_0,y_0)$ is the vertical line $x=x_0$, which meets $E$ again at $(x_0,-y_0)$.
So $-(x,y)=(x,-y)$: \emph{to add $P$ and $Q$, draw the line, take the third point, and reflect it in the $x$-axis}. In coordinates, let $P=(x_1,y_1)$ and
$Q=(x_2,y_2)$ with $Q\ne-P$, and let $\lambda$ be the slope of the line $PQ$, or of the tangent at $P$ if $P=Q$:
\[
\lambda=\frac{y_2-y_1}{x_2-x_1}\quad(x_1\ne x_2),\qquad\lambda=\frac{3x_1^2+a}{2y_1}\quad(P=Q).
\]
Substituting $y=y_1+\lambda(x-x_1)$ into $y^2=f(x)$ gives a cubic in $x$ with leading coefficient $1$ and roots $x_1$, $x_2$ and $x_3$, the abscissa of $P*Q$. The
coefficient of $x^2$ is $-\lambda^2$ on one side and $-(x_1+x_2+x_3)$ on the other, so
\[
P+Q=(x_3,y_3),\qquad x_3=\lambda^2-x_1-x_2,\qquad y_3=-\bigl(y_1+\lambda(x_3-x_1)\bigr).
\]
For Bachet's point $P=(3,5)$ on $y^2=x^3-2$: $\lambda=\frac{27}{10}$, $x_3=\frac{729}{100}-6=\frac{129}{100}$ and $y_3=-\bigl(5-\frac{27}{10}\cdot\frac{171}{100}\bigr)
=-\frac{383}{1000}$. Indeed $\bigl(\frac{383}{1000}\bigr)^2=\frac{146689}{10^6}=\bigl(\frac{129}{100}\bigr)^3-2$.

\begin{remark}
Bachet's point generates infinitely many rational points: $(3,5)$ has infinite order. (By the theorem of Nagell and Lutz, proved in Tome~12, Chapter~12, a
point of finite order on $y^2=x^3+ax+b$ with $a,b\in\mathbb Z$ has integer coordinates; $2P$ does not, so neither $2P$ nor $P$ has finite order.) Two
theorems frame the story. Mordell proved in 1922 that $E(\mathbb Q)$ is always finitely generated, as Poincar\'e had asked in 1901; Tome~12 proves it in
Chapter~12. Mazur proved in 1977 that its finite part is one of exactly fifteen groups, $\mathbb Z/n$ for $1\le n\le10$ or $n=12$ and
$\mathbb Z/2\times\mathbb Z/2m$ for $1\le m\le4$; this is far beyond our reach. The number of copies of $\mathbb Z$, the \emph{rank}, is the subject of a famous open
problem, the conjecture of Birch and Swinnerton-Dyer.
\end{remark}

\section{Nine flexes and twelve lines}

\begin{think}
On the curve $y^2=x^3-x+1$ of Figure~\ref{fig:realcubics}, find by eye the points where the curve changes from bending one way to bending the other. How many
are there, counting $O$? Are they collinear?
\end{think}

When $O$ is a flex, a point $P$ is a flex exactly when the tangent at $P$ meets $E$ three times at $P$, that is, $P*P=P$, that is, $P+P+P=O$: the flexes are
the points of order dividing $3$, the \emph{$3$-torsion} $E[3]$. We count them.

\begin{theorem}\label{thm:nineflexes}
Let $E:y^2=x^3+ax+b$ be smooth over an algebraically closed field of characteristic $0$. Then $E$ has exactly nine flexes, they form a subgroup
$E[3]\cong(\mathbb Z/3)^2$, and three distinct flexes are collinear if and only if they form a line of the affine plane $(\mathbb Z/3)^2$. In particular the
line through two flexes contains a third one, and the nine flexes lie three by three on twelve lines.
\end{theorem}

\begin{proof}
Let $f(x)=x^3+ax+b$. An affine point $P=(x,y)$ with $y=0$ has $-P=P$, so $2P=O$ and $3P=P\ne O$: it is not a flex. For $y\ne0$, $P$ is a flex if and only if
$2P=-P$, which means that $2P$ and $P$ have the same abscissa. By the duplication formula this is $\frac{f'(x)^2}{4f(x)}-2x=x$, that is,
\[
\psi(x)=12xf(x)-f'(x)^2=3x^4+6ax^2+12bx-a^2=0.
\]
Since $f''(x)=6x$,
\[
\psi'(x)=12f(x)+12xf'(x)-2f'(x)f''(x)=12f(x).
\] If $\psi$ had a double root $x_0$, then $f(x_0)=0$, and $\psi(x_0)=0$ would give
$f'(x_0)=0$: a double root of $f$, which smoothness excludes. So $\psi$ has four distinct roots, none of them a root of $f$, and each gives two points
$(x_0,\pm y_0)$ with $y_0\ne0$. With $O$, this is nine flexes.

The flexes are the solutions of $3P=O$, a subgroup of the abelian group $E$. A group of order $9$ in which every element has order $1$ or $3$ is not cyclic,
so it is $\mathbb Z/3\times\mathbb Z/3$. For distinct flexes $P,Q$, the third point $R$ of the line $PQ$ satisfies $P+Q+R=O$, so $R=-P-Q$ is a flex. It is
different from $P$: $R=P$ would give $Q=-2P=P$, since $3P=O$. Likewise $R\ne Q$. In the plane $(\mathbb Z/3)^2$, which is a vector space over $\mathbb F_3$,
the line through $P$ and $Q$ is $\{P,Q,P+2(Q-P)\}$, and $P+2(Q-P)=2Q-P=-P-Q$ since $3=0$. So the lines through two flexes are exactly the affine lines of
$\mathbb F_3^2$. There are $\binom92/\binom32=12$ of them.
\end{proof}

This configuration of nine points and twelve lines is named after Otto Hesse, the first to study its properties; the nine flexes themselves go back to
Maclaurin (see Artebani and Dolgachev in the sources). It is the affine plane of order $3$, which you met in Chapter~4 as the finite plane $\mathbb P^2(\mathbb F_3)$
minus a line. It cannot be drawn with straight lines in the real plane, as we will see in a moment; Figure~\ref{fig:hesse} draws four of its lines as circles.

\begin{figure}[ht]
\centering
\fig{ch08-hesse}
\caption{The Hesse configuration: nine points, twelve lines of three points each, any two points on exactly one line. Eight lines are straight; the four
other ones are drawn as circles.}
\label{fig:hesse}
\end{figure}

\begin{theorem}\label{thm:realflexes}
A smooth real cubic $y^2=x^3+ax+b$ has exactly three real flexes, and they are collinear.
\end{theorem}

\begin{proof}
The real flexes form the subgroup $E(\mathbb R)\cap E[3]$ of $E[3]\cong(\mathbb Z/3)^2$, so there are $1$, $3$ or $9$ of them. There are not $9$: the nine
flexes are not all on one line (a line meets $E$ in at most three points), and every line through two of them contains a third, which the theorem of
Sylvester and Gallai (Exercise~\ref{ex:sylvestergallai}) forbids for a finite set of real points. (The theorem is stated in the affine plane, but a real projective
transformation sending to infinity a real line that misses the nine points brings us there, and it preserves collinearity.) There are more than one: let $e$ be the largest real root of
$f$. Then $\psi(e)=-f'(e)^2<0$, while $\psi(x)\to+\infty$ as $x\to+\infty$, so $\psi$ has a root $x_1>e$, where $f(x_1)>0$. This gives two real flexes
$(x_1,\pm\sqrt{f(x_1)})$, and with $O$ we have three. They are collinear: the vertical line $x=x_1$ passes through $O$.
\end{proof}

\begin{remark}
Every smooth complex cubic has exactly nine flexes, whether or not we start from a Weierstrass equation: they are its intersections with the
\emph{Hessian curve} $\det\bigl(\partial^2F/\partial x_i\partial x_j\bigr)=0$, another cubic, and one shows that all nine intersections are simple (Brieskorn
and Kn\"orrer, or Kirwan, in the sources). In particular every smooth complex cubic has a flex, and Proposition~\ref{prop:weierstrass} applies to it.
\end{remark}

\section{Cubics over finite fields}

Over a finite field $\mathbb F_p$, $p$ an odd prime, the group $E(\mathbb F_p)$ is finite, and counting its points is a game with quadratic residues.
Recall the Legendre symbol: $\bigl(\frac cp\bigr)$ is $1$ if $c$ is a nonzero square mod $p$, $-1$ if it is not a square, and $0$ if $p\mid c$.

\begin{proposition}\label{prop:countEp}
For $E:y^2=f(x)$ over $\mathbb F_p$,
\[
\# E(\mathbb F_p)=p+1+\sum_{x\in\mathbb F_p}\Bigl(\frac{f(x)}p\Bigr).
\]
\end{proposition}

\begin{proof}
The equation $y^2=c$ has $1+\bigl(\frac cp\bigr)$ solutions $y\in\mathbb F_p$: two if $c$ is a nonzero square, one if $c=0$, none otherwise. Summing over $x$
gives $p+\sum_x\bigl(\frac{f(x)}p\bigr)$ affine points, and the point $O$ adds one.
\end{proof}

The sum is a sum of $p$ signs. If they behaved like coin flips it would typically be of size about $\sqrt p$, and that is the truth: Hasse proved in 1933
that $\bigl|\#E(\mathbb F_p)-p-1\bigr|\le2\sqrt p$ (Tome~12, Chapter~13). Figure~\ref{fig:finitecubic} shows the $99$ affine points of $y^2=x^3-x+1$ over
$\mathbb F_{101}$, scattered like stars, but symmetric under $y\mapsto-y$. Sometimes the signs cancel exactly.

\begin{proposition}\label{prop:supersingular}
If $p\equiv3\pmod4$, the curve $y^2=x^3-x$ has exactly $p+1$ points over $\mathbb F_p$.
\end{proposition}

\begin{proof}
Here $f(-x)=-f(x)$, and $-1$ is not a square mod $p$ (Tome~1), so $\bigl(\frac{f(-x)}p\bigr)=-\bigl(\frac{f(x)}p\bigr)$. The terms for $x$ and $-x$ cancel,
the term for $x=0$ is $0$, and the sum in Proposition~\ref{prop:countEp} vanishes.
\end{proof}

\begin{figure}[ht]
\centering
\fig{ch08-finite}
\caption{The $99$ affine points of $y^2=x^3-x+1$ over $\mathbb F_{101}$; with $O$, $E(\mathbb F_{101})$ has $100$ points, and $|100-102|=2\le2\sqrt{101}$.}
\label{fig:finitecubic}
\end{figure}

\begin{remark}
In 1985 Neal Koblitz and Victor Miller independently proposed to use the groups $E(\mathbb F_p)$ for cryptography: the chord-and-tangent rule is easy to
compute, while undoing it (finding $n$ from $P$ and $nP$) seems very hard. Elliptic curve cryptography now protects a large part of the traffic on the
internet. The companion notebook implements the group law over $\mathbb F_p$.
\end{remark}

\section{Exercises}

\begin{exercise}[$\star$]\label{ex:noneutral}
Show that $*$ has no neutral element: there is no $N\in E$ with $N*P=P$ for all $P\in E$.
\end{exercise}

\begin{exercise}[$\star$]\label{ex:collinear}
Let $O$ be a flex. Show that $-P=P*O$, and that $P+Q+R=O$ if and only if $P,Q,R$ are the three points of a line with $E$, counted with multiplicity.
\end{exercise}

\begin{exercise}[$\star$]\label{ex:bachet}
On $y^2=x^3-2$, compute $3P$ for $P=(3,5)$, and check that it is on the curve.
\end{exercise}

\begin{exercise}[$\star$]\label{ex:twotorsion}
Show that the points of order $2$ of $y^2=f(x)$ are the $(e,0)$ with $f(e)=0$, and that over $\mathbb C$ the points with $2P=O$ form a group isomorphic to
$(\mathbb Z/2)^2$. When are all of them real?
\end{exercise}

\begin{exercise}[$\star\star$]\label{ex:torsionsix}
Show that $P=(2,3)$ has order $6$ on $y^2=x^3+1$, and list its multiples.
\end{exercise}

\begin{exercise}[$\star$]\label{ex:realcomponents}
Show that the real points of $y^2=f(x)$ form two pieces if $f$ has three real roots, and one piece if it has one.
\end{exercise}

\begin{exercise}[$\star\star$]\label{ex:fermatflexes}
Find the nine flexes of the Fermat cubic $x^3+y^3+z^3=0$ in $\mathbb P^2(\mathbb C)$, and check that the line through two of them contains a third.
\end{exercise}

\begin{exercise}[$\star\star$]\label{ex:flexesxcube}
Find the flexes of $y^2=x^3+1$. Which of them are real? Check Theorem~\ref{thm:realflexes}.
\end{exercise}

\begin{exercise}[$\star$]\label{ex:countfive}
List the points of $y^2=x^3+x+1$ over $\mathbb F_5$, and show that $E(\mathbb F_5)$ is cyclic. Which points are flexes?
\end{exercise}

\begin{exercise}[$\star\star$]\label{ex:xcubeplusone}
Show that if $p\equiv2\pmod3$, then $y^2=x^3+1$ has exactly $p+1$ points over $\mathbb F_p$.
\end{exercise}

\begin{exercise}[$\star\star$]\label{ex:changeorigin}
Let $O$ and $O_1$ be two points of $E$, with $O$ a flex, and write $+$ and $\oplus$ for the group laws with origins $O$ and $O_1$. Show that
$P\oplus Q=P+Q-O_1$. So the groups are isomorphic, by $P\mapsto P-O_1$.
\end{exercise}

\begin{exercise}[$\star\star\star$]\label{ex:sixonconic}
Let $O$ be a flex, and $P_1,\dots,P_6$ six points of $E$ on a conic $C$. Assume that $A=P_1*P_2$, $B=P_3*P_4$ and $C'=P_5*P_6$ are such that the nine points $P_1,\dots,P_6,A,B,C'$ are distinct and $C'\notin C$. Show that $P_1+\dots+P_6=O$.
\end{exercise}

\section{Hints}

\begin{hint}{ex:noneutral}
If $N*P=P$, the line through $N$ and $P$ meets $E$ at $P$ twice.
\end{hint}
\begin{hint}{ex:collinear}
Use $O*O=O$ in Proposition~\ref{prop:groupbasic}.
\end{hint}
\begin{hint}{ex:bachet}
Add $P$ and $2P=\bigl(\frac{129}{100},-\frac{383}{1000}\bigr)$ with the chord formula.
\end{hint}
\begin{hint}{ex:twotorsion}
$2P=O$ means $P=-P$.
\end{hint}
\begin{hint}{ex:torsionsix}
Compute $2P$ with the tangent, then $3P=P+2P$.
\end{hint}
\begin{hint}{ex:realcomponents}
Look at the sign of $f$ on the intervals between its roots.
\end{hint}
\begin{hint}{ex:fermatflexes}
Try the points with a zero coordinate. For collinearity, use determinants.
\end{hint}
\begin{hint}{ex:flexesxcube}
Here $\psi(x)=3x^4+12x=3x(x^3+4)$.
\end{hint}
\begin{hint}{ex:countfive}
The squares mod $5$ are $0,1,4$. A group of order $9$ is cyclic if it has an element of order $9$.
\end{hint}
\begin{hint}{ex:xcubeplusone}
When $\gcd(3,p-1)=1$, cubing is a bijection of $\mathbb F_p$.
\end{hint}
\begin{hint}{ex:changeorigin}
Compute $P\oplus Q=O_1*(P*Q)$ using $P*Q=-(P+Q)$ and Exercise~\ref{ex:collinear}.
\end{hint}
\begin{hint}{ex:sixonconic}
The cubic $C\cup AB$ passes through eight of the nine points where $E$ meets $P_1P_2\cup P_3P_4\cup P_5P_6$.
\end{hint}

\section{Solutions}

\begin{solution}{ex:noneutral}
\idea A neutral element for $*$ would make every line through it tangent somewhere.
\writeup Suppose $N*P=P$ for all $P$. For $P\ne N$, the line $NP$ then meets $E$ in $N$, $P$, $P$, so by Lemma~\ref{lem:tangentdouble} it is the tangent at
$P$. Now look at the lines through $N$ over $\overline K$. At most one of them (the tangent at $N$) meets $E$ only at $N$; each of the infinitely many others
contains a point $P\ne N$ of $E$ and is the tangent at $P$. So there are infinitely many points of $E$ whose tangent passes through $N$. But these points lie on
the polar conic $n_0F_x+n_1F_y+n_2F_z=0$ of $N=[n]$, which is not zero and has no common component with the irreducible cubic $E$ (solution of
Exercise~\ref{ex:polarcurve}); by B\'ezout there are at most $3\cdot2=6$ of them. Contradiction.
\end{solution}

\begin{solution}{ex:collinear}
\idea With a flex as origin, the tangent at $O$ contributes nothing.
\writeup By Proposition~\ref{prop:groupbasic}, $-P=P*(O*O)=P*O$. If $P,Q,R$ are the points of a line, then $P*Q=R$ and $P+Q=O*R=-R$, so $P+Q+R=O$. Conversely,
if $P+Q+R=O$, let $R'=P*Q$; then $P+Q=O*R'=-R'$, so $-R'+R=O$ and $R=R'$: the points $P,Q,R$ are the three points of a line.
\end{solution}

\begin{solution}{ex:bachet}
\idea Same recipe, with a chord instead of a tangent.
\writeup With $P=(3,5)$ and $2P=\bigl(\frac{129}{100},-\frac{383}{1000}\bigr)$, the slope is
$\lambda=\frac{-\frac{383}{1000}-5}{\frac{129}{100}-3}=\frac{-5383/1000}{-171/100}=\frac{5383}{1710}$. Then
$x_3=\lambda^2-3-\frac{129}{100}=\frac{164323}{29241}$ and $y_3=-\bigl(5+\lambda(x_3-3)\bigr)=-\frac{66234835}{5000211}$, so
\[
3P=\Bigl(\frac{164323}{29241},\ -\frac{66234835}{5000211}\Bigr),\qquad 29241=171^2,\quad 5000211=171^3.
\]
One checks that $y_3^2=x_3^3-2$ (the companion notebook does it in exact arithmetic). The denominators grow fast: this is the typical behaviour of a point of
infinite order.
\end{solution}

\begin{solution}{ex:twotorsion}
\idea Points of order $2$ are their own negatives.
\writeup $2P=O$ means $P=-P$. For $P=(x,y)$ affine, $-P=(x,-y)$, so $y=0$ and $f(x)=0$. Over $\mathbb C$, $f$ has three distinct roots $e_1,e_2,e_3$, so $E[2]$
consists of $O$ and the three points $(e_i,0)$: a group of order $4$ in which every element has order at most $2$, hence $(\mathbb Z/2)^2$. (As a check, the line
$y=0$ contains all three, so $(e_1,0)+(e_2,0)+(e_3,0)=O$.) All of them are real exactly when $f$ has three real roots.
\end{solution}

\begin{solution}{ex:torsionsix}
\idea Compute multiples until you come back to $O$.
\writeup At $P=(2,3)$: $\lambda=\frac{3\cdot4}{6}=2$, $x=4-4=0$, $y=-(3+2(0-2))=1$, so $2P=(0,1)$. Then $P+2P$: $\lambda=\frac{1-3}{0-2}=1$, $x=1-2-0=-1$,
$y=-(3+1\cdot(-1-2))=0$, so $3P=(-1,0)$, a point of order $2$. Hence $6P=O$ and $P$ has order $6$, with $4P=-2P=(0,-1)$ and $5P=-P=(2,-3)$. The multiples are
$O,(2,3),(0,1),(-1,0),(0,-1),(2,-3)$.
\end{solution}

\begin{solution}{ex:realcomponents}
\idea A real point needs $f(x)\ge0$.
\writeup The real points have $f(x)=y^2\ge0$. If $f$ has three real roots $e_1<e_2<e_3$, then $f\ge0$ exactly on $[e_1,e_2]\cup[e_3,\infty)$. Over $[e_1,e_2]$ the
two graphs $y=\pm\sqrt{f(x)}$ join at the ends into a closed oval; over $[e_3,\infty)$ they form an unbounded branch, which closes up through $O$ in the
projective plane. With one real root $e$, $f\ge0$ exactly on $[e,\infty)$, and there is only the unbounded branch.
\end{solution}

\begin{solution}{ex:fermatflexes}
\idea The coordinate lines cut the Fermat cubic in three points each, and those are the flexes.
\writeup Let $\omega=e^{2\pi i/3}$. On the line $z=0$ the curve gives $x^3=-y^3$, so the points $[1:-\omega^k:0]$, $k=0,1,2$; likewise $[0:1:-\omega^k]$ and
$[-\omega^k:0:1]$: nine points. At $P=[1:-1:0]$ the gradient $(3x^2,3y^2,3z^2)$ is $(3,3,0)$, so the tangent is $x+y=0$; on it, $F=(-y)^3+y^3+z^3=z^3$, a triple root
at $z=0$, so $P$ is a flex. Multiplying a coordinate by a cube root of unity, or permuting the coordinates, preserves $F$ and moves $[1:-1:0]$ to any of the nine points, so all nine
are flexes. Collinearity:
$[1:-1:0]$, $[0:1:-1]$, $[-1:0:1]$ satisfy $x+y+z=0$, and in general the determinant of $[1:-\omega^a:0]$, $[0:1:-\omega^b]$, $[-\omega^c:0:1]$ is
$1-\omega^{a+b+c}$, which vanishes when $a+b+c\equiv0\pmod3$. So the line through two flexes on different coordinate lines contains a third, and the three
flexes on one coordinate line are collinear by construction.
\end{solution}

\begin{solution}{ex:flexesxcube}
\idea Solve $\psi=0$ and keep the roots where $f\ne0$.
\writeup With $a=0$, $b=1$: $\psi=3x^4+12x=3x(x^3+4)$. The roots are $x=0$ and the three cube roots of $-4$. At $x=0$, $y^2=1$: the flexes $(0,\pm1)$, real. At
$x=-\sqrt[3]4$, $y^2=-4+1=-3$: $y=\pm i\sqrt3$, not real; at the two complex cube roots $x=-\sqrt[3]4\,\omega^{\pm1}$ also $y^2=-3$. So the nine flexes are $O$, $(0,\pm1)$
and six non-real points, and the three real ones, $O$, $(0,1)$, $(0,-1)$, lie on the line $x=0$. In Exercise~\ref{ex:torsionsix}, $(0,1)=2P$ indeed has order $3$.
\end{solution}

\begin{solution}{ex:countfive}
\idea Tabulate $f(x)$ and compare with the squares.
\writeup $f(0)=1$, $f(1)=3$, $f(2)=11\equiv1$, $f(3)=31\equiv1$, $f(4)=69\equiv4$. The value $3$ is not a square; the others are. So the affine points are
$(0,\pm1)$, $(2,\pm1)$, $(3,\pm1)$, $(4,\pm2)$, and with $O$ there are $9$. For $Q=(2,1)$: $\lambda=\frac{3\cdot4+1}{2}=\frac{13}2\equiv\frac32\equiv3\cdot3=9\equiv4$,
$x=16-4=12\equiv2$, $y=-(1+4\cdot0)=-1$: so $2Q=(2,-1)=-Q$ and $Q$ has order $3$, a flex, as is $-Q$. For $P=(0,1)$: $\lambda=\frac{1}{2}\equiv3$, $x=9\equiv4$,
$y=-(1+3\cdot4)=-13\equiv2$, so $2P=(4,2)$; then $3P=P+2P$ with $\lambda=\frac{2-1}{4-0}=\frac14\equiv4$, $x=16-0-4=12\equiv2$, $y=-(1+4\cdot2)=-9\equiv1$, so
$3P=(2,1)=Q\ne O$. The order of $P$ divides $9$ and is not $1$ or $3$, so it is $9$: the group is cyclic, generated by $P$. Its flexes are $O$, $(2,1)$, $(2,-1)$.
\end{solution}

\begin{solution}{ex:xcubeplusone}
\idea Change variables $u=x^3$.
\writeup If $p\equiv2\pmod3$, then $3\nmid p-1$, so $x\mapsto x^3$ is injective on the cyclic group $\mathbb F_p^\times$ (its kernel has order $\gcd(3,p-1)=1$), hence
bijective on $\mathbb F_p$. So the number of affine solutions of $y^2=x^3+1$ equals that of $y^2=u+1$, which is $p$: one $u$ for each $y$. With $O$, $p+1$ points.
\end{solution}

\begin{solution}{ex:changeorigin}
\idea Every $*$ can be rewritten with $+$, because three collinear points add to $O$.
\writeup By Exercise~\ref{ex:collinear}, $X*Y=-X-Y$ for all $X,Y$. So $P\oplus Q=O_1*(P*Q)=-O_1-(P*Q)=-O_1+P+Q$. The map $\phi(P)=P-O_1$ then satisfies
$\phi(P\oplus Q)=P+Q-2O_1=\phi(P)+\phi(Q)$, and it is a bijection, so it is an isomorphism from $(E,\oplus)$ to $(E,+)$.
\end{solution}

\begin{solution}{ex:sixonconic}
\idea Cayley--Bacharach again, with the conic playing the role of two of the three lines.
\writeup The cubic $L=P_1P_2\cup P_3P_4\cup P_5P_6$ meets $E$ in the nine distinct points $P_1,\dots,P_6,A,B,C'$. The cubic $C\cup AB$ passes through
$P_1,\dots,P_6$ (on $C$) and $A,B$: eight of them. By Theorem~\ref{thm:cayleybacharach} it passes through $C'$, and since $C'\notin C$, $C'$ is on the line $AB$.
So $A$, $B$, $C'$ are collinear, and $A+B+C'=O$ by Exercise~\ref{ex:collinear}. Also $P_1+P_2=-A$, $P_3+P_4=-B$ and $P_5+P_6=-C'$. Adding,
$P_1+\dots+P_6=-(A+B+C')=O$.
\end{solution}

\section*{Sources and further reading}

\begin{itemize}
\item D.~Perrin, \emph{Autour de l'\'equation de Bachet}, lecture notes, Universit\'e Paris-Saclay (Bachet's tangent method of 1621 and the point $\bigl(\frac{129}{100},\frac{383}{1000}\bigr)$
of $y^2+2=x^3$).
\item I.~Newton, \emph{Enumeratio linearum tertii ordinis}, appendix to \emph{Opticks}, London, 1704; the description of Newton's manuscripts at the
Cambridge Digital Library (72 species, projections of the five divergent parabolas).
\item M.~Artebani and I.~Dolgachev, \emph{The Hesse pencil of plane cubic curves}, L'Enseignement Math\'ematique 55 (2009), 235--273 (the Hesse
configuration, its history and its symmetries).
\item J.~H.~Silverman and J.~Tate, \emph{Rational Points on Elliptic Curves}, 2nd ed., Springer, 2015 (the group law, Nagell--Lutz, Mordell's theorem and
curves over finite fields, at the level of this book).
\item H.~Poincar\'e, \emph{Sur les propri\'et\'es arithm\'etiques des courbes alg\'ebriques}, J.\ Math.\ Pures Appl.\ (5) 7 (1901); L.~J.~Mordell, \emph{On the
rational solutions of the indeterminate equations of the third and fourth degrees}, Proc.\ Cambridge Philos.\ Soc.\ 21 (1922); B.~Mazur, \emph{Modular
curves and the Eisenstein ideal}, Publ.\ Math.\ IH\'ES 47 (1977).
\item N.~Koblitz, \emph{Elliptic curve cryptosystems}, Math.\ Comp.\ 48 (1987); V.~Miller, \emph{Use of elliptic curves in cryptography}, CRYPTO '85 (1986).
\item E.~Brieskorn and H.~Kn\"orrer, \emph{Plane Algebraic Curves}, Birkh\"auser, 1986; F.~Kirwan, \emph{Complex Algebraic Curves}, Cambridge University Press,
1992 (the Hessian and the nine flexes).
\end{itemize}
